\documentclass[11pt]{article}
\usepackage[margin=1in]{geometry}
\usepackage{amsmath,amssymb}
\usepackage{enumitem}
\usepackage[T1]{fontenc}
\usepackage[utf8]{inputenc}
\usepackage{lmodern}
\setlength{\parindent}{0pt}
\setlength{\parskip}{0.5em}
\begin{document}

\section*{IB Physics HL: D4 Problem Collection}
Paraphrased induction questions from the available 2025 exam pool.
\begin{itemize}[leftmargin=1.25em]
\item
\textbf{Source.} May 2025 HL Paper 1A TZ2, Question 28\\
\textit{Original format: multiple choice.}\\
A horizontal conducting ring is perpendicular to a uniform vertical magnetic field. The ring is rotated through $180^\circ$ about a horizontal axis in $2.0\,\text{s}$, inducing an average emf of $2.0\,\mu\text{V}$. The area of the ring is $4.0\,\text{cm}^2$. Determine the magnetic field strength.
\item
\textbf{Source.} May 2025 HL Paper 1A TZ2, Question 29\\
\textit{Original format: multiple choice.}\\
A rectangular coil rotates at constant angular speed in a uniform magnetic field. A graph of emf against time is given, along with labelled axes of the coil. Determine which specified rotation axis and angle are consistent with the emf-time graph.
\item
\textbf{Source.} May 2025 HL Paper 1A TZ2, Question 30\\
\textit{Original format: multiple choice.}\\
A magnet attached to a spring oscillates above a coil with period $T$, and the current in the coil is plotted against time. The oscillation amplitude is then reduced. Determine the new current-time graph.
\item
\textbf{Source.} May 2025 HL Paper 2 TZ3, Question 8\\
\textit{Note.} Mixed D2/D4 question; included here because parts (c) and (d) are directly on induction and induced emf.\\
In the same rod-and-wire setup, a circular coil of radius $2.0\,\text{cm}$ with $20$ turns is placed with its centre $10.0\,\text{cm}$ from rod $R$. The current in $R$ is initially constant at $2.0\,\text{A}$.

\begin{enumerate}[label=(\alph*), leftmargin=1.5em]
\setcounter{enumi}{2}
\item
\begin{enumerate}[label=(\roman*), leftmargin=1.5em]
\item Explain why no current is induced in the coil while the current in $R$ remains constant. [2]
\item The current in $R$ increases uniformly from $2.0\,\text{A}$ to $10.0\,\text{A}$ in $0.5\,\text{s}$. Calculate the induced emf in the coil. [3]
\item Deduce the direction of the induced current. [2]
\end{enumerate}
\item The coil is then rotated so that an emf is induced. An emf-time graph is given. The rotation frequency is doubled. Draw the new emf-time graph. [2]
\end{enumerate}
\item
\textbf{Source.} November 2025 HL Paper 2, Question 4\\
\textit{Note.} Mixed D3/D4 question; included here because the entire setup is a motional-induction problem and explicitly invokes Lenz's law.\\
A metal rod moves to the right along rails in a uniform magnetic field into the page.

\begin{enumerate}[label=(\alph*), leftmargin=1.5em]
\item The rod experiences a magnetic force of $0.084\,\text{N}$.
\begin{enumerate}[label=(\roman*), leftmargin=1.5em]
\item Show that the current in the rod is about $0.1\,\text{A}$. [1]
\item State the direction of the magnetic force on an electron in the rod. [1]
\item Explain why the magnetic field magnitude differs on the two sides of the moving rod. [3]
\item Outline how Lenz\'s law applies. [2]
\end{enumerate}
\item Using the current found above, determine the force per unit length between the rails and give its SI base units. [3]
\end{enumerate}
\end{itemize}

\end{document}
